% =====================================================================
%  Цель и задание лабораторной работы.
% =====================================================================

\textbf{Цель:} на основе практических результатов лабораторной работы
№3.1 сформулировать многократно применимую методику оценки безопасности
кода и проверить её работоспособность, сравнив по ней собственную
реализацию с альтернативной, написанной иначе.

\vspace{1em}

\textbf{Задание}:
\begin{enumerate}
  \item На основе полученных результатов предыдущей части лабораторной
  работы создать методику оценки безопасности кода.
  \item По созданной методике сравнить предоставленный алгоритм со своим
  из предыдущей части лабораторной работы.
  \item Оформить отчёт о проделанной работе.
\end{enumerate}

\textbf{Постановка задачи.} Роль «предоставленного алгоритма» (пункт 4
инструментальных средств задания) играет отдельная, самостоятельно
написанная для этой работы программа \texttt{Lab3.2-ProvidedAlgorithm} --
она реализует ровно ту же функциональность, что и раздел «In-memory
data» из лабораторной работы №3.1 (создание/обновление/удаление
конфиденциальных и неконфиденциальных данных в оперативной памяти без
сохранения на диск), но написана так, как эту задачу решил бы
разработчик, не знакомый с выводами лабораторной работы №3.1: без
валидации входных данных средствами фреймворка, с обратимым
«шифрованием» на основе XOR с зашитым в код ключом вместо
AES-256-GCM, с постоянно растущим журналом отладки, в который
по ошибке попадают конфиденциальные значения в открытом виде, и без
синхронизации доступа к общим структурам данных. Такая программа
воспроизводит реальные, часто встречающиеся ошибки, а не гипотетические,
и потому пригодна как контрольный пример для проверки создаваемой
методики: методика имеет смысл только тогда, когда она способна отличить
заведомо небезопасный код от кода, для которого эти вопросы уже были
осознанно проработаны.
